% answers_short.tex —— 2009 年·普通高等学校招生考试·安徽卷·文科 · 简短答案速查
% 由 tools/make_short_answers.py 自 answers.tex 的 \gksolhead/\gkanswer 提取，
% 并与 2009/2009_数学真题_安徽文_答案要点.md 逐题比对后生成，勿手改（改源 MD 后重跑）。

\begin{gknotice}
\textbf{使用说明}：本节只给最终答案便于快速核对，完整推导见解析版。答案为编写者独立推导，\textbf{非官方标准答案}。
\end{gknotice}

\gkpart{选择题}
\begin{center}\begin{tabular}{ccccccccccc}
\toprule
题号 & 1 & 2 & 3 & 4 & 5 & 6 & 7 & 8 & 9 & 10 \\
\midrule
答案 & D & B & C & A & B & B & A & C & D & A \\
\bottomrule
\end{tabular}\end{center}

\gkpart{填空题}
\begin{itemize}[leftmargin=2.6em]
  \item[\textbf{11.}] $(0,-1,0)$。
  \item[\textbf{12.}] $127$。
  \item[\textbf{13.}] $\frac34$。
  \item[\textbf{14.}] $\frac43$。
  \item[\textbf{15.}] ①④⑤。
\end{itemize}

\gkpart{解答题}
\begin{itemize}[leftmargin=2.6em]
  \item[\textbf{16.}] （1）$\dfrac{\sqrt3}{3}$；（2）$3\sqrt2$。
  \item[\textbf{17.}] （1）茎叶图见解析版；（2）优点见下；（3）$A$ 的平均亩产高于 $B$，$B$ 的亩产更稳定。
  \item[\textbf{18.}] （1）$a=\sqrt3$，$b=\sqrt2$；（2）$y^2=-4x$，是抛物线。
  \item[\textbf{19.}] （1）$a_n=4n$，$b_n=2^{1-n}$；（2）证明见解析。
  \item[\textbf{20.}] （1）证明见解析；（2）$2\sqrt2$。
  \item[\textbf{21.}] （1）见解析正文；（2）$\left[2-3\ln2,\ \mathrm e^2-5-2\mathrm e^{-2}\right]$。
\end{itemize}
